% ------------------------------------------------------------------
% Acknowledgements
% ------------------------------------------------------------------

% §1 — Institutional Support
% This work was partially funded by the Fundação para a Ciência e a Tecnologia (FCT), I.P., and developed at the Department of Physics of the University of Coimbra.

% §1b — Personal arc
This work grows out of a journey that began when, as a teenager in Brazil, I observed the planets through a telescope I had built with my own hands---deciding then to become a physicist. I started my studies at Brazil (UFMG) in 2010, spent time in Coimbra in 2011, and in Germany 2014, on this track specizlized in precision optics, working on systems ranging from Alzheimer's disease detectors to optical calibration of satellites (Sentinel-5). That path brought me back to Coimbra, more than a decade later, to complete the Master's in Astrophysics and Space Instrumentation (MAIE) that I had first glimpsed in 2011. Along that journey I witnessed the growing severity of wildfires in both Brazil and Portugal, and chose to dedicate this thesis to a problem that connects the two countries where I have built my life---and the planet I admire. To all who were part of this path---and who allowed me to move from someone who observes other planets to someone capable of conceiving missions to observe and protect the Earth---my most sincere gratitude.

% §2 — Supervisors
Among all, special recognition is owed to my supervisors, whose scientific guidance, intellectual rigour, and constant availability were decisive for the completion of this work.

To Professor Dr. Rui Miguel Curado da Silva, whose expertise in space systems made it possible to transform an initial idea---``Portugal faces wildfires; how can space instrumentation help?''---into a structured mission, connected to wildfire combat groups (ADAI) and to CubeSat fire-detection projects (FADO\footnote{F.A.D.O -- Fire Analysis \& Development Observatory, Concept Report CubeSat 2024, Team CSP-7-FADO.}), in which I participated and from whose colleagues I learned a great deal. His teaching in the space instrumentation course was essential for building the knowledge required to start this work, and his ability to connect people and institutions within the Portuguese space ecosystem was instrumental in defining the directions of this thesis and the future paths for fire-detection sensor development.

To Professor Dr. Carlos Xavier Paes Lidas, whose expertise in wildfire management and combat technologies---namely fire-spread modelling and the use of satellite imagery by firefighting authorities in Portugal---was fundamental in defining what the mission had to fulfil in the context of fire detection in Portugal. Part of the knowledge and data that informed this thesis was derived from two projects connected to Professor Carlos: the project IMFire---Intelligent Management of Wildfires\footnote{This work was carried out under the project IMFire---Intelligent Management of Wildfires, ref.\ PCIF/SSI/0151/2018, partially funded by national funds through the Ministry of Science, Technology and Higher Education, Portugal.}, and the project ForestSphere\footnote{This research was partially funded under the Protocol between Fundação para a Ciência e a Tecnologia and the Agency for Administrative Modernization, I.P.\ (AMA), through the project ForestSphere, ref.\ 2025.01496.DT4ST.}. It was also he who introduced me to Professor João Luís Ruivo Carvalho Paulo and the artificial-intelligence-based fire detection work within the FireSys project.

To Professor Dr. Alan Kardec de Almeida Júnior, whose knowledge of orbital mechanics was essential for advancing from basic mission requirements to high-fidelity models. He taught me to use GMAT, the mission-analysis tool developed at agencies such as NASA, and transmitted to me the fundamentals of celestial mechanics, perturbation forces, atmospheric drag, orbital stability, station-keeping, and the simulation tools required---foundations that constitute the simulation infrastructure presented here. I am equally grateful for his careful guidance in the writing and revision of this work.

To all three, my most sincere gratitude for the support, encouragement, and trust throughout this journey.

% §3 — Project Collaborators
I am equally grateful to Professor João Luís Ruivo Carvalho Paulo for the opportunity to work on the FireSys\footnote{This dissertation work was partially supported by the project FireSys -- ``Intelligent Decision Support System for Management of Wildfires'' with reference of operation 21378, operation code COMPETE2030-FEDER-02230700 on Balcão dos Fundos, funded by the Fundo Europeu de Desenvolvimento Regional (FEDER), through the Innovation and Digital Transition Program (COMPETE 2030) of Portugal 2030.} project on satellite-image-based fire detection using AI. From that project emerged the \emph{Watcher--Zoom} concept presented in this work.
I thank colleagues PhD\,(c) Dina Jahanianfard, who shared valuable insights on satellite image processing related to Sentinel-2 and the demonstration of the \emph{Watcher--Zoom} concept, and PhD\,(c) Matheus Fernandes Kovaleski, who collaborated on the code development.
I thank Dr.\ Timothée Vaillant for his suggestions regarding the defence presentation.
I also wish to thank the MAIE coordinators: Prof.\ Dr.\ Margarida Maria Lopes da Silva Camarinha, whose teaching in the Celestial Mechanics course provided the mathematical foundations and formalism present in this work, and Prof.\ Dr.\ Alexandre C. M. Correia for his support in navigating the decisions that helped define the thesis topic during the Master's and in bringing this work to completion.

% §4 — Scientific and Professional Formation
It is also important to acknowledge the experiences that built the foundations of this work before I arrived in Coimbra, starting in Brazil at UFMG. To Professor Dr. Ado Jório Vasconcelos, for the scientific training in optics at LabNS that constitutes my technical base and for believing in me at the start of my career, and to Professor Dr. Leandro Malard and colleagues at LabNS for their contributions to that formative in a unique environment. To Professor Renato Las Casas at the Observatório da Serra da Piedade, who taught me how to share the sky with others, and to Bernardo Riedel, who taught me the art of building telescopes. To Dr.\ Luiz Etrusco and colleagues at InventVision, for their support during my studies and in the development of advanced cameras. To Dr. Ivo Vieria and colleagues at LusoSpace, for showing me in practice how space systems are built.

% §5 — Colleagues and Friends
I wish to thank all the friends and colleagues who were with me at the University of Coimbra or who were part of this journey to the completion of the Master's, and who at the right moments helped me to relax, enjoy nature, and reflect on ideas.
To Flávia Amaral, astrophysicist and friend who can see vast distances in the universe that I can only begin to imagine, and who has accompanied this journey since 2010 at UFMG. To Pedro Pessoa, mathematician and friend whose mastery of geometry in many dimensions accompanied and enriched our studies in general relativity.
To my Master's friends though this journey ---Cláudia, Letícia, Snock, Aline, Guilherme, Carlos, João Werneck, Akel, Céci, Willian, Ronald, Rita, José Rodrigues, Fred, Maria João, Ricardo, Mariana Encarnação, Inês de Castro, Maria Campino, Fernanda, Joana, Cátia, Hortence, Diogo, Marcos, Angela, and Rubens Silva, Thiago Célio, Thiago Duarte, Filomeno, Eduardo Valadares, Ubirajará Agero, Paula, Renan, Alexandre, Emerson, Fabiano, Marcia, Hudson, Caciano, Lucas, Aroldo, Luiz, Wescley, Camila, Marina, Paola, Bruno, Erlan, André, Rafael, Luis, Kenji, André, Beto.
A special thank-you to Lara and Fig, for showing me that my technical skills in optical instrumentation, combined with space systems, could have a real impact on climate challenges---revealing a meaningful professional path, with consequences for the lives of many people and beings on this planet, and influencing the choice of the topic of this thesis.

% §6 — Family
Finally, I wish to express my deepest gratitude to my family.
To my father Joel and my mother Marisa, who supported me unconditionally at every stage---from the education they provided in Brazil to the decision to study abroad---and who, even at a distance, were never absent.
To my sisters Sabrina, Mariana, and Marina, for their constant encouragement and for always reminding me of what matters.
To all members of the family who contributed, throughout my life, with support, guidance, and the example that sustained effort is the surest way to accompanying our dreams.
This work is also yours.
